% !TeX program = pdflatex
% =============================================================================
% Lektionsplan (Diskussionsgrundlage, NICHT das Arbeitsblatt selbst)
% Magnetfelder -- UZH Praktikum I, Sara Romer Urban, KS im Lee, HS2026
% Klasse: 4fh
%
% Stand 01.09.2026: komplett neu aus arbeitsblatt5-magnetfelder-teil1.tex
% und aufgabenblatt4-magnetfelder-teil1.tex abgeleitet -- die vorherige
% Fassung dieser Datei war veraltet (Ring/Helmholtz/Ringspule und
% Permeabilitaet sind auf Loris' Wunsch nicht mehr Teil dieser
% Doppellektion; die tatsaechliche Abfolge im Arbeitsblatt weicht auch
% sonst deutlich von der alten Planung ab). Jede Tabellenzeile entspricht
% jetzt direkt einem Abschnitt im Arbeitsblatt, in derselben Reihenfolge.
% Neu: pro Phase eine Spalte "Fragen" (was im Unterrichtsgespraech zu
% fragen ist) und "Kernaussagen" (was am Ende jeder Phase sitzen muss).
% Ausserdem klargestellt (Stand 01.09.2026, zweite Korrektur): Das
% Baukasten-Experiment direkt nach Oersted ist die SCHUELERAKTIVITAET
% (Partnerarbeit) und laeuft NUR mit kleinen Kompassen, OHNE Eisenspaene.
% Die beiden Eisenfeilspaene-Demonstrationen (gerader Leiter auf weisser
% Flaeche; Spule auf Plexiglas-Halterung) sind separate, kuerzere
% PLENUM-Demos von vorne, die je direkt VOR der zugehoerigen Formel
% gezeigt werden -- als praezisere Bestaetigung dessen, was die SuS im
% Baukasten bereits selbst mit den Kompassen beobachtet haben.
% =============================================================================
\documentclass[11pt,a4paper]{article}

\usepackage[T1]{fontenc}
\usepackage[utf8]{inputenc}
\usepackage[ngerman]{babel}
\usepackage{lmodern}
\usepackage[a4paper,landscape,margin=1.5cm]{geometry}
\usepackage{array,booktabs,tabularx,xltabular}
\usepackage{enumitem}
\usepackage{xcolor}
\usepackage{amsmath}
\usepackage{siunitx}
\sisetup{per-mode=fraction,fraction-command=\frac}
\usepackage{url}
\Urlmuskip=0mu plus 1mu
\def\UrlBreaks{\do\/\do-\do_}

\definecolor{titleblue}{HTML}{183A6B}
\definecolor{accentblue}{HTML}{2E6F9E}
\definecolor{groupteal}{HTML}{1B7F72}
\definecolor{darkgray}{HTML}{4A4A4A}
\definecolor{warnred}{HTML}{9E2E2E}

\setlength{\parindent}{0pt}
\setlength{\parskip}{0.5em}
\setlist[itemize]{leftmargin=1.1em,itemsep=0.15em,topsep=0.15em}

\newcommand{\Methode}[1]{{\small\color{groupteal}\textbf{#1}}}

\usepackage{titlesec}
\titleformat{\section}{\Large\bfseries\color{titleblue}}{\thesection}{0.6em}{}
\titlespacing*{\section}{0pt}{1.2em}{0.5em}

% Spaltentypen für die Tabelle
\newcolumntype{Y}[1]{>{\raggedright\arraybackslash\small}p{#1}}

\begin{document}
\sloppy

\begin{center}
{\Huge\bfseries\color{titleblue}Lektionsplan: Magnetfelder}\\[0.4em]
{\large Doppellektion (2$\times$45 Min.) -- Diskussionsgrundlage}
\end{center}
\vspace{0.5em}

\noindent\textbf{Klasse:} 4fh (01.09.2026, nur 4fh) \hfill \textbf{Lehrperson:} Loris Keller

\vspace{0.3em}
\noindent\textcolor{darkgray}{Dieses Dokument ist \textbf{keine} Reinschrift des Arbeitsblatts, sondern
die Planungsgrundlage dazu -- 1:1 aus dem fertigen
\texttt{arbeitsblatt5-magnetfelder-teil1.tex} abgeleitet, in derselben
Reihenfolge. Für jede Phase: Inhalt/Methode, was Sie im
Unterrichtsgespräch fragen sollten, und die Kernaussage, die am Ende
sitzen muss.}

% =========================================================================
\section*{1. Abgrenzung und Lernziele}

\textbf{Im Fokus:} korrektes Einzeichnen von Magnetfeldern
(Permanentmagnete, Regeln, Vergleich mit dem elektrischen Feld,
Erdmagnetfeld), das Magnetfeld eines geraden stromdurchflossenen Leiters
(Ørsted-Effekt, Rechte-Faust-Regel, Formel) und das Magnetfeld einer Spule
(experimentell mit Eisenfeilspänen untersucht, Rechte-Faust-Regel, Formel,
Windungsdichte).

\textbf{Bereits bekannt (Werkstatt Magnetismus, nur Repetition):} Pole,
Anziehung/Abstossung, kein magnetischer Monopol, Elementarmagnete-Modell,
Weiss'sche Bezirke, Ferro-/Dia-/Paramagnetismus als Begriff.

\textbf{Bewusst nicht Teil dieser Einheit:} Induktion/Faraday, Lorentzkraft
(eigene spätere Doppellektionen), Ring/Helmholtz-Spule/Ringspule und
magnetische Permeabilität (auf Loris' ausdrücklichen Wunsch aus dieser
Doppellektion entfernt -- die Spulenformel wird direkt eingeführt, ohne den
Umweg über Ring $\to$ zwei Ringe).

\begin{itemize}[leftmargin=1.3cm,itemsep=0.3em]
  \item[LZ1] Magnetfelder korrekt einzeichnen -- Regeln kennen und
  anwenden, den strukturellen Unterschied zum elektrischen Feld benennen
  (kein magnetischer Monopol), das Erdmagnetfeld einordnen.
  \item[LZ2] den Ørsted-Effekt und das Magnetfeld eines geraden,
  stromdurchflossenen Leiters beschreiben (Rechte-Faust-Regel) und die
  Formel $B=\frac{\mu_0}{2\pi}\cdot\frac{I}{r}$ (magnetische Flussdichte)
  anwenden.
  \item[LZ3] das Magnetfeld einer Spule beschreiben (Rechte-Faust-Regel,
  experimentell mit Eisenfeilspänen untersucht) und die Formel
  $B=\mu_0\cdot\frac{N}{l}\cdot I=\mu_0\cdot n\cdot I$ (Windungsdichte $n$)
  anwenden.
\end{itemize}

% =========================================================================
\section*{2. Aufbau der Doppellektion}

\begin{xltabular}{\textwidth}{@{}Y{2.1cm} Y{6.6cm} Y{6.6cm} Y{6.6cm} p{1.1cm}@{}}
\toprule
\textbf{Phase} & \textbf{Inhalt \& Methode} & \textbf{Fragen} & \textbf{Kernaussagen} & \textbf{Zeit} \\
\midrule
\endhead

Repetition &
\textbf{Magnete und Kompassnadel.} Bild: Kompassnadel zwischen zwei
Stabmagneten (N/S einander zugewandt). \Methode{Einzelarbeit, dann
Vergleich zu zweit}. &
\begin{itemize}
  \item Wie richtet sich die Nadel aus, und warum genau so?
  \item Was wissen Sie aus der Werkstatt noch über Pole und
  Elementarmagnete?
\end{itemize} &
\begin{itemize}
  \item Die Kompassnadel folgt den Feldlinien; ihr Nordpol zeigt in
  Feldrichtung (N $\to$ S).
  \item Zwischen ungleichnamigen Polen ist das Feld dicht und stark
  (Anziehung) -- die Nadel schlägt dort deutlich aus.
\end{itemize} &
6' \\[0.6em]

Vergleich &
\textbf{Punktladungsfeld vs. Stabmagnetfeld.} Zwei Bilder nebeneinander
(elektrischer Dipol / Stabmagnet), beide sehen ähnlich aus.
\Methode{Stillarbeit, dann Plenum}. &
\begin{itemize}
  \item Was fällt beim Vergleich der Bilder auf?
  \item Wo liegt der Unterschied, obwohl die Form ähnlich aussieht?
\end{itemize} &
\begin{itemize}
  \item Elektrische Ladungen können einzeln existieren (Monopol möglich).
  \item Beim Magneten gibt es das nicht: Nord- und Südpol sind
  untrennbar -- \textbf{kein magnetischer Monopol}.
\end{itemize} &
6' \\[0.6em]

Theorie &
\textbf{Regeln zum Zeichnen von Magnetfeldern} (5 Regeln), danach ein
korrekt gezeichnetes Beispiel (Stabmagnetfeld) zur Kontrolle.
\Methode{Lehrervortrag, geleitet}. &
\begin{itemize}
  \item Wo beginnen/enden Feldlinien -- aussen und innen des Magneten?
  \item Was bedeutet eine dichtere Ballung von Feldlinien?
\end{itemize} &
\begin{itemize}
  \item Aussen: N $\to$ S. Innen: S $\to$ N, jede Linie geschlossen.
  \item Feldlinien kreuzen sich nie; Richtung = Kompassnadel-Nordpol;
  Dichte = Stärke.
\end{itemize} &
6' \\[0.6em]

Zeichenübung &
\textbf{Zwei Stabmagnete, Nordpole zueinander.} SuS zeichnen die Feldlinien
selbst in eine TikZ-Vorlage (nur die beiden Magnete vorgegeben). Danach
\textbf{Fehlersuche}: eine absichtlich falsche Lösung mit sich
kreuzenden Feldlinien zwischen den Nordpolen wird gezeigt.
\Methode{Einzelarbeit + Plenumsdiskussion}. &
\begin{itemize}
  \item Was passiert zwischen zwei Nordpolen?
  \item (Fehlersuche) Was stimmt an dieser Lösung nicht -- welche Regel
  wird verletzt?
\end{itemize} &
\begin{itemize}
  \item Gleichnamige Pole stossen sich ab; Feldlinien weichen seitlich
  aus, kein Übergang, neutraler Punkt in der Mitte.
  \item Feldlinien dürfen sich nie kreuzen (Fehler in der gezeigten
  Lösung).
\end{itemize} &
10' \\[0.6em]

Theorie &
\textbf{Erdmagnetfeld.} Bild + Hinweisbox geografischer vs. magnetischer
Pol. \Methode{Lehrervortrag, geleitet}. &
\begin{itemize}
  \item Warum zeigt der Kompass nach Norden, wenn Feldlinien doch vom
  magnetischen Nordpol \emph{weg} verlaufen?
\end{itemize} &
\begin{itemize}
  \item Geografischer Nordpol $\neq$ magnetischer Nordpol.
  \item Am geografischen Nordpol liegt der \textbf{magnetische Südpol}
  des Erdfelds.
\end{itemize} &
4' \\[0.6em]

Überleitung &
\textbf{Die Geschichte von Ørsted} (kurz, mit Bild). \Methode{Erzählimpuls}. &
\begin{itemize}
  \item Was könnte eine stromdurchflossene Leitung mit einem Magnetfeld
  zu tun haben?
\end{itemize} &
\begin{itemize}
  \item 1820, Zufallsentdeckung: Strom erzeugt selbst ein Magnetfeld --
  erste Verbindung von Elektrizität und Magnetismus.
\end{itemize} &
3' \\[0.6em]

Baukasten-Experiment &
\textbf{SuS experimentieren selbst: Stromkreis mit Stableiter und Spule,
nur kleine Kompasse (keine Eisenspäne).} Feld an beiden Bauteilen mit
Kompassen erkunden, danach Stromrichtung umpolen und erneut beobachten.
\Methode{Partnerarbeit}. &
\begin{itemize}
  \item Wie richten sich die Kompasse um den Leiter aus, wie in/um die
  Spule?
  \item Was verändert sich an \emph{beiden} Feldern bei Stromumkehr?
\end{itemize} &
\begin{itemize}
  \item Leiter: Kompasse stellen sich tangential zu gedachten Kreisen.
  Spule: innen einheitlich, an den Enden wie ein Stabmagnet, seitlich
  schwächer.
  \item Beide Feldrichtungen kehren sich bei Stromumkehr vollständig um;
  Nord- und Südpol der Spule vertauschen sich.
\end{itemize} &
12' \\[0.6em]

Experiment 1 &
\textbf{Plenum-Demo: gerader Leiter auf weisser Fläche mit
Eisenfeilspänen.} Draht senkrecht durch die Fläche, Strom ein-/
ausschalten, Muster mit Eisenspänen sichtbar machen -- \textbf{direkt vor}
der Formel, als präzisere Bestätigung dessen, was die SuS eben selbst mit
Kompassen im Baukasten beobachtet haben. \Methode{Demonstration, Plenum}. &
\begin{itemize}
  \item Welches Muster bilden die Eisenspäne um den Draht?
  \item Passt das zu dem, was Sie eben mit den Kompassen gesehen haben?
\end{itemize} &
\begin{itemize}
  \item Konzentrische Kreise um den Leiter, bestätigt die eigene
  Kompass-Beobachtung im Detail.
\end{itemize} &
5' \\[0.6em]

Theorie &
\textbf{Formel des geraden Leiters.} Rechte-Faust-Regel (Bild), dann
$B=\frac{\mu_0}{2\pi}\cdot\frac{I}{r}$ -- SuS ergänzen Formel und Einheit
selbst im Antwortfeld, bevor die Lösung gezeigt wird.
\Methode{Lehrervortrag, geleitet, mit Schülerbeteiligung}. &
\begin{itemize}
  \item Welche Grössen könnten $B$ beeinflussen -- und wie (direkt oder
  umgekehrt proportional)?
  \item Welche Einheit könnte $B$ haben?
\end{itemize} &
\begin{itemize}
  \item $B=\frac{\mu_0}{2\pi}\cdot\frac{I}{r}$; $B\propto I$,
  $B\propto\frac1r$.
  \item Einheit: Tesla (T) $=\,\si{\newton\per\ampere\per\meter}$.
\end{itemize} &
8' \\[0.6em]

Experiment 2 &
\textbf{Plenum-Demo: Spule auf Plexiglas-Halterung mit
Eisenfeilspänen.} Feld der Spule mit Eisenspänen sichtbar machen (innen,
an den Enden, seitlich) -- \textbf{direkt vor} der Formel, als präzisere
Bestätigung dessen, was die SuS eben selbst mit Kompassen beobachtet
haben. \Methode{Demonstration, Plenum}. &
\begin{itemize}
  \item Wie unterscheidet sich das Muster von dem des geraden Leiters?
  \item Passt das zu dem, was Sie eben mit den Kompassen gesehen haben?
\end{itemize} &
\begin{itemize}
  \item Im Inneren: starkes, homogenes Feld entlang der Spulenachse.
  \item An den Enden: deutliche Pole wie bei einem Stabmagneten; seitlich
  deutlich schwächer.
\end{itemize} &
5' \\[0.6em]

Theorie &
\textbf{Formel der Spule.} Rechte-Faust-Regel für die Spule (Bild,
\emph{andere} Regel als beim Leiter!), dann
$B=\mu_0\cdot\frac{N}{l}\cdot I$ -- SuS ergänzen Formel und Einheit selbst,
danach Windungsdichte $n=\frac{N}{l}$ einführen.
\Methode{Lehrervortrag, geleitet, mit Schülerbeteiligung}. &
\begin{itemize}
  \item Was passiert mit dem Feld bei mehr Windungen auf gleicher Länge?
  \item Worin unterscheidet sich die Rechte-Faust-Regel hier von jener
  beim geraden Leiter?
\end{itemize} &
\begin{itemize}
  \item $B=\mu_0\cdot\frac{N}{l}\cdot I=\mu_0\cdot n\cdot I$ mit
  Windungsdichte $n=\frac{N}{l}$.
  \item Leiter: Daumen = Stromrichtung. Spule: Finger = Stromrichtung,
  Daumen = Nordpol -- nicht verwechseln.
\end{itemize} &
8' \\
\bottomrule
\end{xltabular}

\textcolor{darkgray}{\small Summe: 73' + 17' Puffer (Fragen, Übergänge,
Materialverteilung bei den beiden Experimenten) = 90'.}

% =========================================================================
\section*{3. Anschlussdokument}

Wie geplant entstanden ist ein separates \textbf{Aufgabenblatt}
(\texttt{aufgabenblatt4-magnetfelder-teil1.tex}, im selben Ordner) für die
Vor-/Nachbereitung: Feldlinien zeichnen, Stabmagnet-Polung bestimmen,
Formelumstellung (gerader Leiter, Spule, jeweils mit Nicht-SI-Einheiten zum
Umrechnen), Verständnisfragen zur Spule, sowie eine Ringspule-Zusatzaufgabe
(bewusst als Vertiefung über den Stoff dieser Doppellektion hinaus, mit
kurzer Erklärung direkt in der Aufgabe). Alle Aufgaben mit
Schwierigkeitsgrad-Angabe und Lösung.

\end{document}
